% =====================================================================
\chapter{Background}
\label{ch:background}
% =====================================================================

This chapter sets out the minimum background needed to follow the rest of the
dissertation. It explains how a transformer language model turns text into internal
vectors, states precisely what we mean by a hallucination, frames detection as a
classification problem, and defines the metrics used to score a detector. A reader
already familiar with transformer internals and detection metrics can skim this chapter
and move on to Chapter~\ref{ch:related}.

\section{Transformer language models and their hidden states}

Modern language models are built from the transformer architecture~\cite{vaswani2017attention}.
The input text is first broken into \emph{tokens} --- sub-word units --- and each token
is mapped to a vector. These vectors then pass through a stack of identical layers. Each
layer mixes information between token positions using a mechanism called
\emph{self-attention}, and then transforms each position independently through a small
feed-forward network. The output of every layer is, again, one vector per token. We
write $\zvec_\ell^{(t)}$ for the vector that layer $\ell$ produces at token position
$t$, and we call it the \emph{activation} or \emph{hidden state} at that layer and
position. A model with $L$ layers therefore produces a whole stack of activations
$\zvec_1^{(t)}, \zvec_2^{(t)}, \dots, \zvec_L^{(t)}$ for each position it processes.

After the final layer, the model multiplies the last activation by an output matrix to
obtain a score for every possible next token, and a softmax turns these scores into a
probability distribution $p_t(\cdot)$ over the vocabulary. The token with the highest
probability is usually emitted, the text grows by one token, and the process repeats.
Two facts about this picture matter for the rest of the dissertation. First, the entire
stack of activations is computed anyway during normal generation, so anything we derive
from it is essentially free. Second, the activations are high-dimensional --- a few
thousand numbers per token for the models we study --- and they are believed to encode
not only \emph{what} the model is saying but also something about \emph{how sure} and
\emph{how grounded} it is. The central bet of internal-state detection is that this
``something'' can be read out.

The dimensionality of a model's hidden state, written $d$, and its number of layers $L$
vary with model size. The models used in this dissertation range from a small diagnostic
model with a few hundred million parameters up to seven-billion-parameter models; the
six-billion-parameter model that anchors our results has $d = 4096$ and $L = 28$ layers.
These sizes are reported in full in Chapter~\ref{ch:method}.

\section{What counts as a hallucination}

The word ``hallucination'' covers more than one phenomenon, and being careful about the
distinctions avoids confusion later. Following the standard
survey~\cite{ji2023survey}, we separate two axes.

The first axis is \emph{intrinsic} versus \emph{extrinsic}. An intrinsic hallucination
contradicts the material the model was given: a summary that reverses what an article
actually said is intrinsic. An extrinsic hallucination adds content that cannot be
checked against the given material --- it may happen to be true or false in the world,
but it is unsupported by the input.

The second axis is \emph{faithfulness} versus \emph{factuality}. Faithfulness asks
whether the output is consistent with the prompt or source document; factuality asks
whether it is consistent with the real world. A model can be faithful but not factual
(it accurately reflects a flawed source) or factual but not faithful (it states a true
fact the source never mentioned). Different benchmarks target different corners of this
space, which is one reason a detector that does well on one benchmark may do poorly on
another.

A third consideration is \emph{granularity}: a hallucination label can be attached to a
whole passage, to a sentence, or to an individual fact. Finer labels are more useful but
harder to obtain. The approach we build on produces labels automatically at the level of
a generated continuation, which sidesteps the cost of human annotation; the price is
that the labels are only as good as the automatic procedure that creates them, a point
we return to in Chapter~\ref{ch:method}.

Throughout the dissertation we use the operational definition adopted by the method we
extend: a hallucination is a coherent but factually incorrect span of generated text,
where ``incorrect'' is judged against an external reference such as an encyclopedia
article or a gold question--answer pair.

\section{Detection as a classification problem}

Once labels exist, detection becomes a supervised classification problem. For each
generated example we form a feature vector $\mathbf{x}$ from the model's internal state,
and we attach a binary label $y$, where $y=0$ marks a faithful generation and $y=1$
marks a hallucinated one. A classifier is then trained to estimate
$P(y = 1 \mid \mathbf{x})$. In this work the classifier is a multi-layer perceptron,
abbreviated MLP: a feed-forward neural network of several layers that maps the feature
vector to a probability. The detector's quality depends almost entirely on whether
$\mathbf{x}$ carries the right information, which is why the design of the features ---
the subject of Chapter~\ref{ch:method} --- is where the real work lies.

A subtle but important point is the difference between two evaluation regimes. If the
classifier is trained and tested on the same kind of text, we measure
\emph{in-distribution} detection. If it is trained on one kind of text (encyclopedia
continuations) and tested on another (question answering, summarisation, dialogue), we
measure \emph{transfer}. Transfer is the harder and more realistic setting, and as
Chapter~\ref{ch:results} shows, the two regimes can give very different numbers for the
same detector. We never average across them.

\section{How detection is measured}

Because a detector outputs a probability, it can be scored either as a ranking or as a
calibrated estimate. We use the following metrics, each defined here at first use.

\textbf{AUROC.} The area under the receiver-operating-characteristic curve, abbreviated
AUROC, measures how well the detector \emph{ranks} hallucinated examples above faithful
ones. It equals the probability that a randomly chosen hallucinated example receives a
higher score than a randomly chosen faithful one. A perfect detector scores $1.0$, and
a detector that guesses scores $0.5$. AUROC is the headline metric in most of the
internal-state literature because it does not depend on a particular decision threshold,
and it is the metric we use for cross-method comparison.

\textbf{Accuracy, precision, recall, F1.} These threshold-dependent metrics summarise
the confusion matrix once a cut-off is chosen. Accuracy is the fraction of correct
decisions; precision is the fraction of flagged examples that really were hallucinated;
recall is the fraction of true hallucinations that were flagged; and F1 is the harmonic
mean of precision and recall. They are useful but can be misleading when the two classes
are imbalanced or when the classifier collapses to predicting a single class, a failure
mode we observe at small data scales.

\textbf{Brier score and ECE.} A detector is well \emph{calibrated} if a stated
probability of, say, $0.8$ is right about $80\%$ of the time. The Brier score is the
mean squared difference between the predicted probability and the true label; the
expected calibration error, abbreviated ECE, bins predictions by confidence and measures
how far the average confidence in each bin is from the observed accuracy. Lower is better
for both. We report these where our measurement files contain them, because a calibrated
probability is more useful in practice than a bare yes/no verdict.

A final caution applies to all of these numbers. When a benchmark is small, a metric
estimated on it is noisy, and a single easy or hard dataset can dominate an average
computed over several benchmarks. We flag both effects explicitly when they occur in
Chapter~\ref{ch:results}, rather than letting a headline average paper over them.
